%!TEX root=../main.tex
\section{Discussion} % (fold)
\label{sec:discussion}

n this work, we presented a novel representation of NBA players’ shot density charts based on Functional Data Analysis (FDA). We defined the smoothed densities of each player’s made and missed shots as functional data over the entire offensive half-court. This functional representation enabled a principal component decomposition, in which a player’s variation from the mean shot chart is parsimoniously described by a common set of principal component functions, weighted by player-specific scores.

We first used these principal component functions to describe the major modes of variation in shooting behavior. The first component, which explained almost $70\%$ of the total variance, primarily captured shooting frequency, distinguishing between players who shoot with high and low volumes. The remaining components, which explained successively less variation, captured increasingly subtle but interpretable features of shot selection.

We then applied k-medoids clustering to the player-specific scores to construct a data-driven taxonomy of player shooting tendencies. The resulting clusters revealed diverse shooting profiles, which we interpreted through average component scores within each group.

Interestingly, the estimated clusters showed only weak agreement with the conventional NBA-defined player positions. This may be because shooting behavior represents only one aspect of a player’s role; play-making, defensive responsibilities, and off-the-ball movement also influence positional classification. Nevertheless, our analysis provides a novel data-driven perspective on how players differ in their shooting behaviors.

We envision a number of potential uses of both our functional shot-chart representation and the subsequent clustering of shooting patterns.
Firstly, the functional representation of a players shot-chart provides a comprehensive characterization of a player's shooting behavior, i.e., the shot chart is not reduced to one or more scalar summaries.
This representation could be used by coaches or front-office staff to monitor changes in a player's shooting patterns between or within seasons, to identify player-specific tendencies which inform specific training regimes, and also to identify players whose characteristics match the needs of the team during scouting and recruitment.
Because the charts are can naturally be presented in a way that is interpretable to basketball fans who may not be familiar with the underlying mathematics and statistics, the representation can be used to tell stories to the general public about player tendencies, and enable historical comparisons between players, enriching the narrative around the sport.
While we used shots made and shots missed, other metrics defined over the offensive half-court (e.g., shooting efficiency), could be represented in the same way for the aforementioned purposes.

There are a number of possible methodological extensions to the modeling framework.
One of the most exiting avenues is to combine the shot-chart representations with other high-throughput data sources that are now routinely collected within the NBA. 
For example, temporal biomechanical data are now being measured by motion-tracking systems such as Hawk-Eye, and the possibility of a multi-modal representation that reflects a players tactical and physical tendencies -- and the interaction between them -- could provide more detailed and specific analyses.
The models could also be applied in other sports, for example, to characterize a player's typical position in soccer.

{\color{red}Add references}

% {\color{red}In-season use -- is a player's pattern changing?}

% As an extension, use the defensive shot chart from \cite{franksCharacterizingSpatialStructure2015}.

% As an extension, use the shots efficiency and frequency and not shots attempted and made.


% \textcolor{red}{
% The practical applications of analyzing basketball shooting patterns using functional data analysis and clustering can significantly enhance various aspects of player development, coaching strategies, and team performance. Here’s a detailed exploration of these applications:}

% \textcolor{red}{
% **1. Player Development:**
%    - **Tailored Training Programs:** By identifying specific shooting patterns and inefficiencies, coaches can design personalized training regimens that focus on improving weak areas. For example, if a player's analysis shows a lower shooting percentage from certain court locations, targeted drills can be implemented to enhance their skills in those areas.
%    - **Biomechanical Adjustments:** Insights from functional data analysis can help players refine their shooting mechanics. Understanding the optimal release angles, velocities, and body positioning can lead to improved shooting accuracy. For instance, adjusting the release height or speed based on biomechanical feedback can enhance performance[1][3].}

% \textcolor{red}{
% **2. Coaching Strategies:**
%    - **Game Preparation:** Coaches can utilize shooting pattern analyses to prepare specific game strategies against opponents. By studying the shooting tendencies of opposing players, coaches can develop defensive schemes that effectively limit high-percentage shots from preferred locations[2][4].
%    - **In-Game Adjustments:** Real-time data analysis during games allows coaches to make informed decisions about substitutions or play calls based on player performance metrics. For example, if a player is consistently missing shots from a particular area, they may be advised to alter their shot selection or take a break to reassess their approach.}

% \textcolor{red}{
% **3. Team Performance Optimization:**
%    - **Shot Selection Analysis:** Teams can analyze collective shooting data to determine the most effective shot selections during games. This includes evaluating the success rates of different shot types (e.g., three-pointers vs. mid-range shots) and making strategic decisions about when to attempt specific shots based on situational context[4][5].
%    - **Enhanced Scouting Reports:** Detailed analyses of shooting patterns allow for more comprehensive scouting reports on both players and teams. This information can be invaluable for identifying potential trade targets or assessing the strengths and weaknesses of upcoming opponents[2][3].}

% \textcolor{red}{
% **4. Fan Engagement and Analytics:**
%    - **Visualizations for Fans:** Shot charts and other visual analytics tools can be used to engage fans by providing deeper insights into player performances and game strategies. This not only enhances the viewing experience but also fosters a greater understanding of the game among fans[2][4].
%    - **Data-Driven Storytelling:** Teams and analysts can leverage shooting data to tell compelling stories about player development, game trends, and historical comparisons, enriching the narrative around the sport.}


% \textcolor{red}{
% 1. Biomechanics of the Basketball Jump Shot: This source discusses how biomechanics impact shooting performance and provides insights into effective training methods that could be informed by data analysis[1].
% 2. A Model-Based Approach to Shot Charts Estimation in Basketball: This paper outlines advanced methodologies for estimating shot patterns and emphasizes the importance of accurate visualizations in basketball analytics[2].
% 3. New Technology Changing NBA Shooting: This article highlights how emerging technologies are revolutionizing player training through detailed biomechanical analysis, which complements functional data analysis approaches[3].
% 4. Shot Distribution in the NBA: This research examines how shot patterns have evolved over time in relation to rule changes like the three-point line, offering valuable context for current shooting trends[4].
% 5. Using Data to Investigate NBA Shot Patterns: This project demonstrates practical applications of data science in analyzing NBA shooting data, providing insights into shooter accuracy relative to distance and defensive pressure[5].}

% Citations:
% [1] https://www.physio-pedia.com/Biomechanics_of_the_Basketball_Jump_Shot
% [2] https://arxiv.org/html/2405.01182v1
% [3] https://www.espn.co.uk/nba/story/_/id/38042314/i-thought-was-very-informative-new-technology-change-nba-shooting-forever
% [4] https://link.springer.com/article/10.1007/s12662-020-00690-7
% [5] https://nycdatascience.com/blog/student-works/using-data-to-investigate-nba-shot/
% [6] https://www.stat.uga.edu/sites/default/files/Bradley_Lecture%20II.pdf
% [7] https://www.ncbi.nlm.nih.gov/pmc/articles/PMC9494817/
% [8] https://www.frontiersin.org/journals/psychology/articles/10.3389/fpsyg.2022.917980/full

% section discussion (end)